\Needspace{10\baselineskip}
\pdfbookmark[0]{Notices. Mémorial du dépôt général de la guerre. 1830.}{anzeigen-memorial-1830}
\begin{center}
\rule{.85\textwidth}{.3pt}\par
Göttingische gelehrte Anzeigen.\quad 29 April 1830.\par
\rule{.85\textwidth}{.3pt}
\end{center}
\hspace*{\parindent}\textit{Mémorial du dépôt général de la guerre, imprimé par ordre du ministre.} Vol.\,I, 1829 (for 1802---1803), 696 pp. Vol.\,III, 1826 (for 1825), 466 pp. Vol.\,IV, 1828 (for 1826), 494 pp. Vol.\,V, 1829 (for 1827---1828), 490 pp., in quarto. With numerous maps and plans. Paris, \textsc{Picquet}.

The military journal named above began in 1802 at the initiative of General \textsc{Andreossy}, then director of the Depot. The aim was gradually to make some of the treasures of this remarkable institution, immeasurably enriched during the Revolutionary Wars, available for the benefit of military science and its auxiliary disciplines, including history, topography, geodesy, statistics and so on. The respite afforded by the peace that had just come was to be put to use for this purpose. The first issues of the journal thus appeared in rapid succession. When war soon broke out again, however, its progress gradually stalled, and with the seventh
